\section{可复现性说明}
\label{sec:reproducibility-zh}

全部内置数据均由源代码和固定随机种子生成，不下载或依赖受版权保护的 1945 年后乐谱。处理后语料保存为 \path{data/processed/project2_v3_main.pt}；其 20,000/2,000/2,000 显式划分和 SHA-256 摘要记录在 \path{results/project2_full_split_summary.json}。该摘要记录随机种子 42、\texttt{smoke=false}，并说明覆盖周期训练未丢弃任何乐谱主体标记。

在 Windows 上，\path{scripts/smoke_project2.ps1} 运行隔离的 CPU 验证流程，\path{scripts/run_project2_full_local.ps1 -Resume} 运行并审计 CUDA 实验矩阵。正式环境使用 Python 3.11.9、PyTorch 2.5.1+cu121、CUDA 和 NVIDIA GeForce RTX 4060 Ti。实测进程内存峰值为 4.637GiB，CUDA 已分配显存峰值为 0.667GiB，未发生显存不足调整。

每个神经训练目录都包含最佳检查点、可恢复的最后检查点，以及记录优化器状态、缩放器状态、随机数状态、轮次历史、最佳轮次、配置哈希、语料哈希、运行时间和资源峰值的摘要。全部神经行均达到 60 轮。在 \texttt{pcset\_only} 训练前曾出现一次 Windows CUDA 驱动启动故障；日志记录的健康检查和恢复运行在不改变数据或设置的情况下完成了该配置。事件保存在 \path{results/project2_full_run_incidents.json}。

规范汇总输出为 \path{results/project2_metrics.csv}、\path{results/project2_constraints.csv} 和 \path{results/project2_generation_examples.json}。完整命令历史、检查点哈希、已完成配置、环境审计和产物门控位于 \path{results/project2_full_run_report.md}。论文表格由规范 CSV 生成，图~\ref{fig:full-results-zh} 的源 CSV 与构建脚本保存在 \path{paper/figures/}。20 样例专家包及其清单位于 \path{expert_eval/project2/}。

实验完成版本记录在完整运行报告中，当前中文稿没有引入新的实验数值。公开归档应包含代码、配置、哈希、指标文件、表格、代表性 MusicXML 和再生成说明；大型张量和检查点可以根据归档平台容量重新生成或单独保存。
